\subsection{SIGNS AND WORDS}

* Let S be a non-empty set, whose elements will be called signs (this terminology being appropriate to the metamathematical application we have in mind). Let \(\mathrm{L}_{0}(\mathrm{S})\) be the free semigroup generated by S; the elements of \(\mathrm{L}_{0}(\mathrm{S})\) are called words and are identified with finite sequences \(A = (s_{i})_{0 \leqslant i \leqslant n}\) of elements of S. The law of composition in \(\mathrm{L}_{0}(\mathrm{S})\) will be written multiplicatively, so that AB is the sequence obtained by juxtaposition of A and B. The empty word \(\phi\) is the identity element of \(\mathrm{L}_{0}(\mathrm{S})\) . We recall that the length \(l(A)\) of a word \(A \in \mathrm{L}_{0}(\mathrm{S})\) is the number of elements in the sequence A; thus \(l(AB) = l(A) + l(B)\) , and the words of length 1 are the signs. Let \(\mathrm{L}(\mathrm{S})\) denote the set of non-empty words in \(\mathrm{L}_{0}(\mathrm{S})\) .

Suppose, moreover, that we are given a mapping \(s \to n(s)\) of \(S\) into the set \(\mathbf{N}\) of integers \(\geqslant 0\) . For each non-empty word \(A = (s_i)_{0 \leqslant i \leqslant k}\) of \(\mathbf{L}(\mathbf{S})\) , we put

\[
n (A) = \sum_ {i = 0} ^ {k} n (s _ {i})
\]

and \(n(\phi) = 0\) . \(n(A)\) is called the weight of \(A\) . Clearly \(n(AB) = n(A) + n(B)\) . If \(A = A'BA''\) , the word \(B\) is said to be a segment of \(A\) (a proper segment if also \(B \neq A\) ). If \(A'\) (resp. \(A''\) ) is empty, \(B\) is said to be an initial (resp. final) segment of \(A\) . If \(l(A') = k\) , then \(B\) is said to begin at the \((k + 1)\) th place.

If \(A = BCDEF\) (where the words \(B, C, D, E, F\) may be empty) the segments \(C\) and \(E\) of \(A\) are said to be disjoint.

